\documentclass[11pt, a4paper]{article}

% --- Sprache und Enkodierung ---
\usepackage[utf8]{inputenc}
\usepackage[T1]{fontenc}
\usepackage[ngerman]{babel}
\usepackage{amsmath}


% --- Seitenränder ---
\usepackage{geometry}
\geometry{
    a4paper,
    left=25mm,
    right=25mm,
    top=25mm,
    bottom=25mm,
    headheight=14pt
}

% --- Typografie & Schriftart ---
\usepackage{microtype}
\usepackage{lmodern}
\renewcommand{\familydefault}{\sfdefault}

% --- Tabellen & Grafiken ---
\usepackage{tabularx}
\usepackage{booktabs}
\usepackage{graphicx}
\usepackage{float}

% --- Hyperlinks ---
\usepackage{hyperref}
\hypersetup{
    colorlinks=true,
    linkcolor=blue,
    urlcolor=blue
}

% --- Kopf- und Fußzeilen ---
\usepackage{fancyhdr}
\pagestyle{fancy}
\fancyhf{}
\fancyhead[L]{\textsl{Good Knight VR}}
\fancyhead[R]{\textsl{Spielanleitung}}
\fancyfoot[C]{\thepage}
\renewcommand{\headrulewidth}{0.4pt}
\renewcommand{\footrulewidth}{0.0pt}

\title{Good Knight VR Game\\ \large Spielanleitung}
\author{Hikmet Gözaydin 85987}
\date{\today}

\begin{document}

\maketitle
\begin{center}
    \vspace{2cm}
    \includegraphics[width=0.4\textwidth]{Images/image.png}
\end{center}
\newpage

\section{Schritt-für-Schritt-Anleitung}

\subsection{1. Aufsitzen (Start zu Fuß)}
Das Spiel beginnt in den Stallungen der Arena, wo Sie sich zu Fuß befinden.
\begin{enumerate}
    \item Gehen Sie auf Ihr Pferd zu.
    \item Strecken Sie Ihre rechte Hand aus und platzieren Sie sie über dem Sattel des Pferdes.
    \item Halten Sie den \textbf{Grip-Button} (mit dem rechten Mittelfinger) gedrückt.
    \item Ihre Kamera-Sicht schnappt augenblicklich in den Sattel ein, und Sie befinden sich in der Reitposition.
\end{enumerate}

\subsection{2. Zum Startgatter reiten}
Suchen Sie die leuchtende Startzone (Ready Zone) auf der Jousting-Bahn:
\begin{itemize}
    \item \textbf{Ungerade Runden (1, 3, 5):} Reiten Sie Ihr Pferd in das vordere Gatter (\textit{Forward Gate}).
    \item \textbf{Gerade Runden (2, 4, 6):} Reiten Sie Ihr Pferd in das hintere Gatter (\textit{Return Gate}).
\end{itemize}
Der Schiedsrichter startet den Countdown erst, wenn Ihr Pferd ruhig innerhalb der Begrenzung des Startgatters steht.

\subsection{3. Bereitschaft signalisieren}
Sobald Sie sich in der korrekten Position im Gatter befinden:
\begin{enumerate}
    \item Drücken Sie die \textbf{Y-Taste} auf Ihrem linken Controller.
    \item Bereiten Sie sich während des 3-sekündigen Countdowns auf den Anritt vor:
    \begin{quote}
        \textbf{3...} BEREITEN SIE IHRE LANZE VOR \\
        \textbf{2...} DECKUNG NEHMEN / SCHILD BEREITEN \\
        \textbf{1...} ANGRIFF!
    \end{quote}
\end{enumerate}

\subsection{4. Der Anritt \& Zusammenstoß}
Nutzen Sie den \textbf{linken Joystick} Ihres Controllers, um das Pferd die Bahn entlangzureiten und zu lenken.
\begin{itemize}
    \item \textbf{Angriff:} Bewegen Sie Ihre physische rechte Hand, um Ihre Lanze präzise auf die Trefferzonen des Gegners (Helm/Kopf oder Torso) auszurichten. Halten Sie dabei die Grip-Taste (Mittelfinger) gedrückt, um die Lanze zu greifen.
    \item \textbf{Verteidigung:} Heben Sie Ihre linke Hand, um Ihren Schild schützend vor die berechnete Lanzenbahn des Gegners zu halten. Halten Sie dabei die Grip-Taste (Mittelfinger) gedrückt, um den Schild zu greifen.
    \item Bei einem erfolgreichen Treffer ertönt ein satter, metallischer Plattenrüstungskollision-Sound.
\end{itemize}

\subsection{5. Zurücksetzen für den nächsten Durchgang}
Bringen Sie Ihr Pferd nach dem Zusammenstoß manuell zum Stehen (indem Sie den linken Joystick zurückziehen). Während der KI-Gegner eigenständig abbremst und einen Spurwechsel zum gegenüberliegenden Startgatter durchführt, müssen Sie Ihr Pferd wenden und zum nächsten aktiven Gatter reiten, um den nächsten Durchgang zu starten.

\section{Steuerung \& Tastenbelegung}
Die folgende Tabelle zeigt die Tastenbelegung für VR-Controller (Meta Quest 3s) und Tastatur:

\vspace{0.5cm}
\noindent
\begin{tabularx}{\textwidth}{l X}
    \toprule
    \textbf{Aktion} & \textbf{VR-Eingabe (Oculus / Index)} \\
    \midrule
    \textbf{Aufsitzen (Mount)} & Grip-Taste (Mittelfinger) mit rechter Hand nahe dem Sattel gedrückt halten. \\
    \textbf{Steuern \& Reiten} & Linker Joystick (zur Richtungs- \& Geschwindigkeitskontrolle). \\
    \textbf{Bereitschaft (Ready)} & Y-Taste (linker VR-Controller). \\
    \textbf{Spiel pausieren} & Menü- / Pause-Taste. \\
    \textbf{Lanze ausrichten} & Physische Bewegung der rechten Hand (VR-Controller). Dabei muss die Grip-Taste (Mittelfinger) gedrückt gehalten werden, um die Lanze festzuhalten. \\
    \textbf{Schild positionieren} & Physische Bewegung der linken Hand (VR-Controller). Dabei muss die Grip-Taste (Mittelfinger) gedrückt gehalten werden, um den Schild festzuhalten. \\
    \bottomrule
\end{tabularx}
\vspace{0.5cm}

\section{Turnierregeln \& Siegbedingungen}
Die Entscheidung über den Turniersieg fällt auf eine von drei Arten:

\begin{enumerate}
    \item \textbf{Absoluter Sieg (Vom Pferd stoßen):} Reduzieren Sie die Lebenspunkte (Health Pool) des Gegners von $100$ auf $0$, um ihn aus dem Sattel zu werfen.
    \item \textbf{Technischer Knockout (Entwaffnung):} Treffen Sie wiederholt den rechten Arm des Gegners, um dessen Waffenarm-Lebenspunkte ($30\,\text{HP}$) auf $0$ zu reduzieren. Dadurch verliert er seine Lanze und das Turnier gilt als gewonnen.
    \item \textbf{Entscheidung nach Punkten:} Überleben beide Kontrahenten alle 6 Durchgänge, gewinnt der Ritter mit den meisten verbleibenden Lebenspunkten. Bei gleichem HP-Stand endet das Turnier unentschieden.
\end{enumerate}

\end{document}
